% GENERATED by CasDesignBasis/tools/design_report.py from design_basis.toml v2.4.11. Do not edit.
\begin{tabularx}{\textwidth}{@{}l r l l L@{}}
\toprule
Symbol & Value & Unit & Source & Reason \\
\midrule
$N_\text{stroke}$ & \num{2080} & pulses & [FD] 3.2.3 & Oil-flow pulses for a full stroke: about 5200 pulses/l times about 0.4 l of oil per stroke. \\
$t_\text{open}$ & \num{30} & s & - & Time for a full opening stroke with the hydraulic pump running. [FD] leaves it open ("ca nn sekunder").\assumed{} \\
$t_\text{close}$ & \num{30} & s & - & Time for a full closing stroke. Equal to the opening stroke until the cylinder areas and the pump flow are known.\assumed{} \\
$p_\text{load}$ & \num{40} & bar & - & Pressure on the driven line while the cylinder moves. Must stay below the PLC end-position limit (CylinderConfigDUT.MaxPressureOpen/Close, default 100 bar).\assumed{} \\
$p_\text{relief}$ & \num{160} & bar & - & Pressure on the driven line when the piston stands against an end stop (or is jammed) and the pump runs: the relief valve setting. Above the PLC end-position limit, so the PLC sees the end position.\assumed{} \\
$p_\text{idle}$ & \num{5} & bar & - & Residual pressure on a line that is not driven.\assumed{} \\
$\tau_p$ & \num{0.3} & s & - & Time constant of the line pressure after a valve or end-stop change.\assumed{} \\
$t_\text{hpu}$ & \num{1} & s & - & Time from the hydraulic pump start until it delivers flow. Shorter than the PLC PumpStartUpTime (5 s), so the cylinders move once the PLC opens a valve.\assumed{} \\
\bottomrule
\end{tabularx}
